\documentclass{article}
\usepackage{style}

\title{LADR, 3.B}

\begin{document}
\maketitle

\section*{Problem 1}
Give an example of a linear map $T$ such that $\dim\nil T=3$ and $\dim\range T=2$.

\subsection*{Solution}
By fundamental theorem of linear maps this is a map from a five-dimensional space, such that three dimensions are mapped to zero, and two dimensions are mapped to the rest of the co-domain.

\vs

Thus $T : \F^5\to\F^2$ where $T(x,y,z,v,w)=(v,w)$.

\section*{Problem 2}
Suppose $V$ is a vector space and $S, T\in\L(V, V)$ are such that
\[\range S\subset\nil T\]
Prove that $(ST)^2=0$

\subsection*{Solution}
Let $v\in V$, and let $w\in V$ such that $STv=w$. Since $w\in\range S$ and $\range S\subset\nil T$, it follows $w\in\nil T$. Thus $STSTv=STw=S0$. Since all linear maps map zero to zero, $S0=0$. Thus for all $v\in V, (ST)^2v=0$, and so $(ST)^2=0$ as desired.

\section*{Problem 3}
Suppose $\cons(v,m)$ is a list of vectors in $V$. Define $T\in\L(\F^m,V)$ by
\[T(\cons(z,m))=\suml(z,v,m)\]

\subsection*{Solution}
(a) What property of $T$ corresponds to $\cons(v,m)$ spanning $V$?

If $\cons(v,m)$ spans $V$ it follows $\range T=V$, i.e. $T$ is surjective.

\vs

(b) What property of $T$ corresponds to $\cons(v,m)$ being linearly independent?

If $\cons(v,m)$ is linearly independent, the only way to write $\suml(z,v,m)=0$ is by setting all of $\cons(z,m)$ to zero. Thus only $T(0)=0$. Since null space equals $\{0\}$ is equivalent to injectivity (lemma 3.16), $T$ is injective.

\section*{Problem 9}
Suppose $T\in\L(V,W)$ is injective and $v_1,\ldots,v_n$ is linearly independent in $V$. Prove that $Tv_1,\ldots,Tv_n$ is linearly independent in $W$.

\subsection*{Solution}
Suppose for contradiction $Tv_1,\ldots,Tv_n$ is not linearly independent in $W$. Then for $\alpha_1,\ldots,\alpha_n\in\F$:
\[\alpha_1Tv_1+\ldots+\alpha_n Tv_n=0\]

such that at least one $\alpha_i\neq 0$. By linearity:
\[T(\alpha_1v_1+\ldots+\alpha_n v_n)=0\]

Since $T$ is injective, $\nil T=\{0\}$, and therefore $\alpha_1v_1+\ldots+\alpha_n v_n=0$. There is at least one $\alpha_i\neq 0$ in this equation. But $v_1,\ldots,v_n$ is linearly independent and we have a contradiction. Thus $Tv_1,\ldots,Tv_n$ is linearly independent in $W$ as desired.

\section*{Problem 20}
Suppose $W$ is finite-dimensional and $T\in\L(V,W)$. Prove that $T$ is injective if and only if there exists $S\in\L(W,V)$ such that $ST$ is the identity map on $V$.

\subsection*{Solution}
Suppose $T$ is injective. It follows $\dim V\leq\dim W$ (see lemma 3.23). Let $v_1,\ldots,v_n$ be a basis of $V$, and let $w_1,\ldots,w_n\in W$ such that $Tv_i=w_i$ for $i=1,\ldots,n$. Because $T$ is injective, $w_1,\ldots,w_n$ is linearly independent (see hw 8, problem 9). Further, $v_i$ is the only distinct vector mapped by $T$ to $w_i$ for all $i$. Extend $w_1,\ldots,w_n$ to $w_1,\ldots,w_n,w_{n+1},\ldots,w_m$ , a basis in $W$.

\vs

Now we can define $S$. Let $Sw_i=v_i$ for $i=1,\ldots,n$, and $Sw_j=0$ for $j=n+1,\ldots,m$. For all other $w\in W$ linearly extend $Sw=a_1 w_1+\ldots+a_n w_n + a_{n+1} w_{n+1} a_m w_m$, where $a_{1\ldots m}\in\F$.

Then for all $u\in V$:
\begin{align*}
STu&=ST(\alpha_1v_1+\ldots+\alpha_n v_n)\\
&=S(\alpha_1w_1+\ldots+\alpha_n w_n)\\
&=\alpha_1v_1+\ldots+\alpha_n v_n=u    
\end{align*}
where $\alpha_1,\ldots,\alpha_n\in\F$. Thus if $T$ is injective, there exists $S\in\L(W,V)$ such that $ST$ is the identity map on $V$ as desired.

\vs

Conversely, suppose there exists $S\in\L(W,V)$ such that $ST$ is the identity map on $V$. Suppose for contradiction $T$ is not injective. Let $u,v\in V$. Since $T$ is not injective, by definition $Tu=Tv$ does not imply $u=v$. Applying $S$ to both sides we get $STu=STv$, and thus $u=v$. We have a contradiction, and thus $T$ is injective as desired.

\section*{Problem 21}
Suppose $V$ is finite-dimensional and $T\in\L(V,W)$. Prove that $T$ is surjective if and only if there exists $S\in\L(W,V)$ such that $TS$ is the identity map on $W$.

\subsection*{Solution}
Suppose $T$ is surjective, i.e. $\range T=W$. By rank-nullity theorem $\dim V\geq\dim W$. Let $w_1,\ldots,w_n$ be a basis of $W$, and let $v_i,\ldots,v_n\in V$ such that $Tv_i=w_i$ for $i=1,\ldots,n$. It follows that $v_1,\ldots,v_n$ is linearly independent in $V$ (see hw7, problem 4). We can now define $S$ as follows: let $Sw_i=v_i$, and linearly extend $S$ to every vector $w\in W$. Then for all $w\in W$:
\begin{align*}
    TSw&=TS(\alpha_1w_1+\ldots+\alpha_n w_n)\\
    &=T(\alpha_1v_1+\ldots+\alpha_n v_n)\\
    &=\alpha_1w_1+\ldots+\alpha_n w_n=w
\end{align*}
for $\alpha_1,\ldots,\alpha_n\in\F$. Thus if $T$ is surjective there exists $S\in\L(W,V)$ such that $TS$ is the identity map on $W$, as desired.

\vs

Conversely, suppose there exists $S\in\L(W,V)$ such that $TS$ is the identity map on $W$. Suppose for contradiction $T$ is not surjective. Then there exists $w\in W$ such that for no $v\in V,\ Tv=w$. However, $TSw=Tv=w$, which is a contradiction. Thus $T$ is surjective, as desired.

\newpage
\section*{Problem 24}
Suppose $W$ is finite-dimensional and $T_1,T_2\in\L(V,W)$. Prove that $\nil T_1\subset \nil T_2$ if and only if there exists $S\in\L(W,W)$ such that $T_2=ST_1$.

\subsection*{Solution}
Suppose $\nil T_1\subset \nil T_2$. Let $v_1,\ldots,v_k\in V$ be a basis of $\nil T_1$. Since $\nil T_1\subset \nil T_2$, $\nil T_1$ is a subspace of $\nil T_2$. Thus we can extend $v_1,\ldots,v_k\in V$ to $v_1,\ldots,v_k, v_{k+1},\ldots,v_m\in V$, a basis of $\nil T_2$. Finally we can further extend it to $v_1,\ldots,v_k, v_{k+1},\ldots,v_m,v_{m+1},\ldots,v_n\in V$, a basis of $V$.

\vs

Observe that $v_i$ for $i>k$ is a basis in $V$ of vectors that $T_1$ maps to $\range T_1$, and $v_i$ for $i>m$ is a basis in $V$ of vectors that $T_2$ maps to $\range T_2$. We can now define $S$ such that $S(T_1v_i)=T_2v_i$ for $i>m$, and $S(T_1v_i)=0$ for $i\leq m$. Further, we linearly extend $S$ to every vector in $W$.

\vs

For all $v\in V$, $T_2v=T_2(\alpha_1v_1+\ldots+\alpha_n v_n)=\alpha_1T_2v_1+\ldots+\alpha_n T_2v_n$ for $\alpha_1,\ldots,\alpha_n\in\F$. Further:
\begin{align*}
    ST_1v&=ST_1(\alpha_1v_1+\ldots+\alpha_n v_n)\\
    &=\alpha_1ST_1v_1+\ldots+\alpha_n ST_1v_n\\
    &=\alpha_1T_2v_1+\ldots+\alpha_n T_2v_n
\end{align*}

Thus $T_2=ST_1$ as desired.

\vs

Conversely, suppose there exists $S\in\L(W,W)$ such that $T_2=ST_1$. Suppose for contradiction $\nil T_1\not\subset \nil T_2$. Let $v\in\nil T_1$ such that $v\not\in\nil T_2$. Thus $T_1v_1=0$, but $T_2v\neq0$. Since linear maps must map $0$ to $0$, $T_2v=ST_1v=S0=0$. Thus we have a contradiction, and therefore $\nil T_1\subset \nil T_2$ as desired.

\section*{Problem 29*}
Suppose $\phi\in\L(V,\F)$. Suppose $u\in V$ is not in $\nil \phi$. Prove that
\[V=\nil\phi\oplus\{au:a\in\F\}\]

\subsection*{Solution}
Since $u\in V$  is not in $\nil\phi$, it follows $\nil\phi\cap\{au:a\in\F\}=\{0\}$. It further follows $\dim\range\phi\neq 0$, and since $\dim \F=1$ then $\dim\range\phi=1$. Let $n=\dim\nil\phi$. By rank-nullity theorem $\dim V=n+\dim\range\phi=n + 1$.

\vs

Let $v_1,\ldots,v_n\in V$ be a basis of $\nil\phi$. By definition of direct sum:
\begin{align*}
&\nil\phi\oplus\{au:a\in\F\}\\
&=\{(\alpha_1 v_1+\ldots+\alpha_n v_n)+au: \alpha_1,\ldots,\alpha_n,a\in\F\}
\end{align*}

Since $\nil\phi\cap\{au:a\in\F\}=\{0\}$ we know $\alpha_1 v_1+\ldots+\alpha_n v_n+au$ are linearly independent in $V$. Further we know that a list of linearly independent vectors in a space of the right length is a basis of that space. Observe that the length of the list above is $n+1$, which we established is the dimension of $V$. Thus $V=\nil\phi\oplus\{au:a\in\F\}$ as desired.

\section*{Problem 30*}
Suppose $\phi_1$ and $\phi_2$ are linear maps from $V$ to $\F$ that have the same null space. Show that there exists a constant $c\in\F$ such that $\phi_1=c\phi_2$.

\subsection*{Solution}


\section*{Problem 31*}
Give an example of two linear maps $T_1$ and $T_2$ from $\R^5$ to $\R^2$ that have the same null space but are such that $T_1$ is not a scalar multiple of $T_2$.

\subsection*{Solution}

Let $(x, y, z, v, w)\in\R^5$. Let $T_1(x, y, z, v, w)=(v+w,0)$, and let $T_2=(0, v+w)$. Both have the same null space but aren't scalar multiples of each other.


\end{document}